% ============================================================
%  Bagian 3 – Pembahasan Metodologi
% ============================================================

\section{Pembahasan Metodologi}

\subsection{Eksplorasi Data (EDA)}

Eksplorasi dimulai dengan memeriksa tipe data, kelengkapan nilai, duplikasi, dan nilai yang tidak valid. Data latih dan data uji diperiksa secara terpisah. Keduanya memiliki nilai hilang pada delapan fitur yang sama. Jumlah nilai hilang paling banyak terdapat pada \texttt{num\_children}, yaitu 3.035 pada data latih dan 747 pada data uji. Temuan ini menjadi dasar untuk menentukan penanganan nilai hilang sebelum pelatihan.

Distribusi target hanya dapat diperiksa pada data latih. Dari 27.373 pemesanan, 18.404 tidak dibatalkan (67,23\%) dan 8.969 dibatalkan (32,77\%). Karena jumlah pemesanan di kedua kelas berbeda, skor F1 makro digunakan agar kinerja pada masing-masing kelas mendapat bobot yang sama.

Eksplorasi juga mencakup sebaran fitur numerik, frekuensi kategori, hubungan antarfitur, dan pola berdasarkan waktu kedatangan. Untuk fitur seperti jeda pemesanan, tarif per malam, dan kanal pemesanan, sebaran pada data latih dapat dibandingkan menurut status pembatalan. Perbandingan ini menunjukkan asosiasi, bukan sebab pembatalan.

% ----------------------------------------------------------
\subsection{Data Cleaning}

\subsubsection{Penanganan Nilai Hilang}
% TODO: Strategi imputasi (mean, median, modus, drop) beserta alasan.

\subsubsection{Penanganan Data Duplikat}
% TODO: Metode deteksi dan penghapusan duplikat.

\subsubsection{Penanganan Data Tidak Valid}
% TODO: Outlier, nilai negatif pada kolom non-negatif, dsb.

% ----------------------------------------------------------
\subsection{Data Transformation}

\subsubsection{Encoding Fitur Kategorikal}
% TODO: One-hot encoding, label encoding — alasan pemilihan per fitur.

\subsubsection{Normalisasi dan Standardisasi Fitur Numerik}
% TODO: Min-max scaling vs.\ Z-score — alasan pemilihan per algoritma.

\subsubsection{Penanganan Ketidakseimbangan Kelas}
% TODO: SMOTE, undersampling, class weight — metode yang dipilih dan alasannya.

% ----------------------------------------------------------
\subsection{Feature Selection}
% TODO: Metode (filter/wrapper/embedded), fitur yang dipilih, alasan.

\subsection{Dimensionality Reduction}
% TODO: PCA atau lainnya — apakah diterapkan, mengapa.

% ----------------------------------------------------------
\subsection{Implementasi Model}

\subsubsection{K-Nearest Neighbors}
% TODO: Deskripsi kelas KNN — atribut, metode, kompleksitas.
% Contoh template:
%
% \paragraph{Kelas \texttt{KNNClassifier}}
% Kelas ini mengimplementasikan algoritma \textit{K-Nearest Neighbors} dari awal (\textit{from scratch}).
% Atribut dan metode utama dijelaskan berikut.
%
% \begin{lstlisting}[caption={Implementasi KNNClassifier}, label={lst:knn}]
% # tempel source code di sini
% \end{lstlisting}

\subsubsection{Decision Tree Learning}
% TODO: ID3 / C4.5 / CART — kelas, metode rekursi, kriteria split.

\subsubsection{Logistic Regression}
% TODO: Gradient descent, fungsi sigmoid, regularisasi.

\subsubsection{Support Vector Machine}
% TODO: Optimasi dual, kernel trick, parameter C dan gamma.

\subsubsection{Bayesian Network (Naive Bayes)}
% TODO: Estimasi prior dan likelihood, Laplace smoothing.

\subsection{Implementasi Pembanding dengan scikit-learn}
% TODO: Konfigurasi parameter scikit-learn yang dipakai sebagai baseline.

\subsection{Penyimpanan dan Pemuatan Model}
% TODO: Format simpan model (pickle, joblib, dsb.) dan cara memuatnya.

